\scanpage{251}

% head, continuation of scan 250
\mspage{4}\unsure{derniers|deuxièmes} mbs des (4)

% upper display with scratched left label
\[H^{\mathbb D}_{\bullet}(X,M),\quad H_{\mathbb D}^{\bullet}(X,M)\ ,\tag*{\deleted{\unsure{(6)}}}\]

% top paragraph following display, eight lines
sont bien, l'un covariant, l'autre contravariant en $X$. Donc il y a, à isom unique près, une seule paire de bifoncteurs (5) satisfaisant (bi)fonctoriellement des isom. (4), \unsure{ssi|si} les \deleted{foncteurs}\inserted{expressions} (6), considérées comme fonctorielles en $X$, transforment quasi-isomorphismes en isomorphismes. Cela équivaut donc à la validité du

% middle, theorem, two vertical strokes at left
\uline{Théorème} (\unsure{?|6}) Soit $f:X'\longrightarrow X$ dans $\mathrm{Diag}$. Si $f$ est un quasi-isom, c'est une $\mathbb D$-équivalence.

% left of theorem, written obliquely rising to the right
% slanted
\margin{Nouvel axiome ?}

% below theorem through middle of lower third
Quand $X$ \unsure{$\to$|est} la cat. ponctuelle, ce n'est \inserted{\unsure{là|l'}} autre que l'axiome Dér 3 (p. 67). \deleted{\unsure{Nous allons}} J'ignore si c'est une conséquence des axiomes Der 1 à der 5 (1). Si ce n'est pas le cas, il faudrait l'ajouter aux axiomes, je crois. Je n'ai pas même prouvé qu'une équiv. de catégories \inserted{$\in\mathrm{Diag}$} (au sens des catégories !) est une $\mathbb D$-équivalence !

% before (7)
Notons la \uline{formule d'adjonction}

% (7), adjunction
\[\mathrm{Hom}(\unsure{\xi|\zeta}\times_{\mathbb D}M,N)\unsure{\xrightarrow{\sim}|\simeq}
\mathrm{Hom}(M,\mathrm{hom}_{\mathbb D}(\xi,N))\tag*{(7)}\]

% between (7) and (8)
Si $\xi$ provient \deleted{\unsure{Hot}} de $X\in\deleted{\unsure{\mathcal C}}\mathrm{Diag}$, \unsure{donc|dès} cela s'écrit

% (8), last line
\[\mathrm{Hom}(p_!p^*M,N)\xrightarrow{\sim}\mathrm{Hom}(M,p_*p^{\unsure{'|*}}N)\ ,\tag*{(8)}\]

% archive sheet number: 251
